\documentclass[fontset=fandol,zihao=-4]{ctexart}
\usepackage[paperwidth=120mm,paperheight=85mm,margin=10mm]{geometry}
\usepackage{graphicx,unicode-math,user-marker}
\setmainfont{lmroman10-regular.otf}
\setmathfont{latinmodern-math.otf}
\pagestyle{empty}
\newsavebox{\formulabox}
\newwrite\slotfile
\input{formula-metrics}
\begin{document}
% Preserve source scale; its depth locates the baseline within the cropped PDF.
\sbox{\formulabox}{\raisebox{-\sourceformuladepth}{%
\includegraphics[width=\sourceformulawidth]{formula.pdf}}}
\immediate\openout\slotfile=slots.csv
\immediate\write\slotfile{formula-1,\the\wd\formulabox,\the\ht\formulabox,\the\dp\formulabox}
\immediate\closeout\slotfile
\noindent 中文排版与原生公式验证。按基线插入：%
\usebox{\formulabox}，原来的分数线和字形随 PDF 一起迁移。

\noindent 数学字体：$E=mc^2$, $\int_0^1 x^2\,\mathrm{d}x=\frac13$。

\noindent \usermarker
\input{user-template}

% Same position, text and fonts on both proof pages. Compare their actual
% rendered pixels, rather than accepting a PDF solely because its text exists.
\clearpage
\noindent 基线对照：A$\frac{x+1}{y}$B。
\clearpage
\noindent 基线对照：A\usebox{\formulabox}B。
\end{document}
